\chapter{Đặt vấn đề}

\section{Lời hứa trong điều phối đi chung xe}

Đi chung xe theo yêu cầu (on-demand ride-pooling) ghép các yêu cầu có thể phục vụ chung vào
cùng một phương tiện. Bộ điều phối phải cân bằng cơ hội nhận khách mới với thời gian chờ,
thời gian trên xe và sức chứa. Một hướng giải tiêu biểu là sinh các nhóm khách khả thi rồi
chọn phương án gán nhóm cho xe bằng quy hoạch nguyên \cite{alonsomora2017}. Khi yêu cầu xuất
hiện liên tục, việc cập nhật kế hoạch là một phần của quá trình vận hành.

Mỗi kế hoạch còn tạo ra thông tin gửi tới hành khách: xe nào đến đón, đón ở đâu và dự kiến
đến điểm trả lúc nào. Chèn một khách mới có thể làm thay đổi thông tin đã công bố cho khách
cũ. Giữ thời gian chờ dưới mười phút chưa đồng nghĩa với giữ một giờ đón đã báo trước đó.
Hai kế hoạch cùng thỏa cửa sổ phục vụ có thể tạo ra những chuỗi cập nhật rất khác nhau đối
với người đang chờ xe.

Chẳng hạn, một hành khách nhận các giờ đón 9:10, 9:13 rồi 9:11. So giờ cuối với giờ đầu cho
chênh lệch một phút; cộng độ lớn hai lần cập nhật cho bốn phút. Đây là hai câu hỏi khác nhau:
khách cuối cùng đến muộn bao nhiêu, và dự báo đã thay đổi bao nhiêu trong quá trình chờ.
Ví dụ chỉ minh họa cách đo; khóa luận không mặc định mỗi giây cập nhật có cùng tác động tới
tâm lý hoặc chi phí của hành khách.

Độ tin cậy của thông tin này đã được nghiên cứu. Fielbaum và Alonso-Mora
\cite{fielbaum2020} đo những thay đổi so với thông báo ban đầu và thử các quy tắc hạn chế
chúng. Các nghiên cứu lựa chọn dịch vụ cũng cho thấy hành khách quan tâm tới thời gian chờ
ngoài dự kiến \cite{alonsogonzalez2020,gerzinic2023}. Vì vậy, phát hiện bất ổn hoặc yêu cầu
giữ lời hứa tự nó không phải đóng góp mới của khóa luận.

Vấn đề cụ thể ở đây là xác định một nghĩa vụ mà bộ điều phối kiểm được trước khi công bố
kế hoạch, đồng thời nói rõ nghĩa vụ ấy còn có ý nghĩa gì khi thời gian đi thay đổi. Nếu xe
chạy chậm hơn, dự báo có thể xấu đi ngay cả khi bộ điều phối giữ nguyên tuyến. Nếu toàn bộ
thay đổi đó bị tính vào một ngân sách cứng, phương án giữ tuyến cũng có thể bị loại. Nếu chỉ
tính phần thay đổi do lựa chọn kế hoạch, hệ thống giữ được ngân sách hành vi nhưng hành
khách vẫn có thể nhìn thấy một độ trễ lớn hơn ngân sách ấy.

\begin{khungy}
Khóa luận nghiên cứu lời hứa dưới dạng một đại lượng được định nghĩa, ghi sổ và kiểm chứng
trong mô hình điều phối. Câu hỏi trung tâm là: hệ thống đang giới hạn thay đổi do quyết định
của mình, hay đang giới hạn toàn bộ thay đổi mà hành khách nhìn thấy?
\end{khungy}

\section{Khoảng trống và câu hỏi nghiên cứu}

Các cơ chế gần nhất bảo vệ những đối tượng khác nhau. Có cơ chế siết thời gian chờ theo dự
báo; có cơ chế giữ độ chùng cho từng hành khách; có cơ chế đặt hạn chót từ thông báo đầu
tiên \cite{fielbaum2020,elmi2026,schulz2026}. Chúng đều có thể kiểm tra khi chọn phương án.
Do đó, khoảng trống được xét không phải là sự thiếu vắng mọi ràng buộc lời hứa trong văn liệu.
Khóa luận tập trung vào quan hệ giữa cơ sở ghi sổ, lịch sử cập nhật và tính khả thi dưới
trôi dạt, trong cùng một hệ thống thực nghiệm.

RideBound lưu lịch sử cam kết theo từng hành khách. Ở mỗi quyết định, hệ thống đối chiếu
ba mốc: thông tin đã công bố, thông tin thu được khi giữ tuyến cũ trên trạng thái hiện tại,
và thông tin của phương án mới. Đối chiếu này xác định phần thay đổi gắn với lựa chọn
phương án ngay tại trạng thái đó. Nó là một phép phân rã trong mô hình, không phải ước
lượng tác động nhân quả tổng thể của chính sách lên giao thông hoặc hành vi người dùng.

Từ cách đặt vấn đề ấy, khóa luận trả lời ba câu hỏi:

\begin{enumerate}[leftmargin=1.2cm,itemsep=5pt,topsep=4pt]
\item \textbf{RQ1 --- Có thể cưỡng chế đại lượng nào?} Một sổ cái theo hành khách cần lưu
những gì để kiểm ngân sách trước khi chấp nhận phương án, và chứng cứ quyết định xác nhận
được điều gì?
\item \textbf{RQ2 --- Chi phí dịch vụ là bao nhiêu?} Khi giữ cùng bộ sinh phương án và
cùng đầu vào, giảm biến động lời hứa đánh đổi bao nhiêu yêu cầu được phục vụ? Kết luận thay
đổi ra sao khi đổi mức ngân sách hoặc cách đo?
\item \textbf{RQ3 --- Trôi dạt và phục hồi thay đổi lời hứa thế nào?} Các cổng theo hành vi,
hạn chót và biến động nhìn thấy khác nhau ở khả năng giữ phương án giữ tuyến ra sao? Khi
phải phục hồi sau vi phạm, số khách được phục vụ còn phản ánh mức giữ lời hứa đến đâu?
\end{enumerate}

RQ1 cần đối chiếu mô hình, mã nguồn và chứng cứ từng quyết định. RQ2 cần so sánh ghép cặp,
vì hai chính sách có thể phục vụ những tập khách khác nhau. RQ3 cần tách trạng thái cổng
cam kết loại phương án khỏi trạng thái phương án đã vi phạm cửa sổ đón hoặc một bất biến
vật lý. Cách tách này ngăn việc quy mọi lần hệ thống dừng cho cơ chế lời hứa.

\section{Mục tiêu và đóng góp}

Mục tiêu của khóa luận là xây dựng một cơ chế cam kết có thể kiểm tra trong điều phối cuốn
chiếu, rồi đánh giá cả lợi ích đo được và chi phí của nó. Phần phương pháp sử dụng tổng biến
thiên của các cập nhật, thay vì chỉ lưu chênh lệch giữa thông báo đầu và cuối. Ngân sách áp
theo từng hành khách và từng chiều cam kết; một mức trung bình thấp không thay thế được
kiểm tra trần của người chịu thay đổi lớn nhất.

Khóa luận có ba nhóm đóng góp, với mức độ khẳng định khác nhau.

\paragraph{Cơ chế vận hành và khả năng kiểm chứng.}
RideBound hiện thực phép đối chiếu ba mốc, sổ cái tích lũy, các khóa cam kết và bộ kiểm tra
phương án. Chứng cứ quyết định ghi ngân sách trước và sau cùng các lý do chấp nhận hoặc
từ chối. Đây là đóng góp về mô hình hóa và hiện thực hệ thống: cùng một runner được gọi từ
các bộ chuyển đổi mô phỏng để giữ thống nhất ngữ nghĩa. Kiểm chứng ở đây xác nhận điều kiện
trong mô hình và dữ liệu đầu vào, không chứng nhận mọi giờ đến thực tế của dịch vụ.

\paragraph{Khung so sánh có giả thiết rõ ràng.}
Khóa luận đặt bốn họ cổng vào cùng một khung: tổng phần do quyết định, hạn chót một chiều,
độ lệch hai chiều từ thông báo đầu, và tổng biến động nhìn thấy. Kết quả lý thuyết chỉ ra
điều kiện mà cổng hành vi giữ được phương án giữ tuyến và các chiều trôi dạt có thể làm
những cổng khác loại phương án đó. Giá trị nằm ở việc phát biểu và kiểm tra ranh giới của
từng nghĩa vụ. Chứng minh là sơ cấp; nó không bảo đảm rằng toàn bộ bài toán luôn có nghiệm
khi cửa sổ đón hoặc những ràng buộc vật lý đã hỏng.

\paragraph{Bằng chứng ghép cặp và kết quả âm.}
Thực nghiệm đo sự đánh đổi giữa biến động lời hứa và tỉ lệ phục vụ, sau đó kiểm tra tác động
của trôi dạt và hai luật phục hồi. Các kết quả về giảm biến động được báo cùng phần nhu cầu
không được phục vụ. Những lần dự đoán sai, cổng bị ép phải miễn trừ, hoặc phép so sánh không
đạt điều kiện đều được giữ lại. Phần đối chiếu thước đo và đường cong đánh đổi với nghiên
cứu trước là tái lập độc lập; nó không được tính thành phát hiện mới về sự tồn tại của
bất ổn hay về việc giữ lời hứa có chi phí.

Ranh giới tính mới được giới hạn ở tổ hợp sổ cái theo hành khách, phụ thuộc lịch sử nhiều
lần cập nhật, ngân sách nhiều chiều và chứng cứ có thể kiểm tra, cùng khung phân tích cơ sở
ghi sổ dưới trôi dạt. Chèn động, cập nhật giờ đến, giới hạn ETA, tái phân công, tính nhất quán
giờ phục vụ, nhiều kế hoạch và trì hoãn cam kết đều có tiền lệ. Khóa luận cũng chưa có bằng
chứng trực tiếp về tăng mức hài lòng, tăng phúc lợi hoặc cải thiện hiệu quả vận hành tổng thể.
Chương 2 đối chiếu các cơ chế gần nhất để giữ ranh giới này cụ thể.

\section{Phạm vi nghiên cứu}

Nhu cầu thực nghiệm được dẫn xuất từ dữ liệu chuyến taxi Manhattan, thông qua bộ dữ liệu
FleetPy \cite{fleetpy2026}. Các ô chính lấy từ năm ngày 14--18/11/2018, với bốn mẫu nhu cầu mỗi ngày. Mỗi ô
dùng 108 yêu cầu; cấu hình đội xe và cửa sổ phục vụ được khóa theo từng thí nghiệm.
Đây là các mẫu nhỏ để so sánh có kiểm soát, không phải phát lại toàn bộ nhu cầu taxi
Manhattan. Mạng dẫn xuất là bao đóng thời gian đi giữa 96 điểm chọn, có tọa độ tổng hợp;
nó không phải bản đồ 96 nút của mạng đường phố thật.

Thời gian đi được xét ở nhiều chế độ: giữ cố định, áp hệ số theo thời gian từ dữ liệu,
hoặc thêm nhiễu có kiểm soát. Dữ liệu hệ số giao thông theo ngày là chuỗi dùng chung cho
mạng, không mô tả từng cung đường và không tạo ra một bộ mô phỏng giao thông vi mô. Cụm
``chỉ giao thông đo thật'' trong chương thực nghiệm chỉ chế độ không thêm nhiễu tổng hợp,
vẫn trong mô hình mạng dẫn xuất này.

Các kết quả được phân tầng theo cách thiết kế thí nghiệm. Panel chính dùng để đánh giá
chi phí và kiểm chứng cơ học. Các panel phát triển dùng để tìm và giải thích trường hợp
phân kỳ. Một phép kiểm tra cuối chia 49 ô hợp lệ thành 15 ô hiệu chỉnh và 34 ô kiểm tra,
với ba khung giờ. Các ô này vẫn dùng những ngày đã xuất hiện trong nghiên cứu, và một số
khung giờ chồng lấn nhu cầu. Tách hiệu chỉnh với kiểm tra giảm việc dùng cùng kết quả để
chọn ngưỡng và đánh giá ngưỡng; nó chưa tạo ra bằng chứng độc lập ở quần thể, ở thành phố
khác hoặc trên dữ liệu triển khai mới.

Toàn bộ so sánh trong khóa luận là so sánh nội bộ giữa những cấu hình được nêu rõ. Cùng
hạt giống giúp phát lại một lần chạy; nó không biến mọi hành khách hoặc mọi quyết định
thành mẫu độc lập. Chương 4 ghi đơn vị phân tích, mẫu số và quy tắc ghép cặp của từng kết
quả. Các con số tuyệt đối không được xếp hạng trực tiếp với nghiên cứu có đội xe lớn hơn,
mạng khác hoặc mô hình lựa chọn khách khác.

\section{Tổ chức khóa luận}

Chương 2 trình bày bài toán điều phối và phân biệt các cơ chế đo, giới hạn và công bố lời
hứa trong văn liệu. Chương 3 xây dựng ba mốc, sổ cái, bốn họ cổng, các giả thiết của kết quả
lý thuyết và kiến trúc kiểm chứng. Chương 4 trả lời ba câu hỏi nghiên cứu bằng kết quả chính
và những kiểm tra thăm dò cần thiết, với giới hạn đặt cạnh bằng chứng. Chương 5 tổng hợp
điều đã xác lập, điều còn chưa được chứng minh và những bước tiếp theo có thể kiểm tra.

